% --------------------------------------------------------------------
%  DIAGRAM NODE STYLES
% --------------------------------------------------------------------
%  The vocabulary the hand-drawn TikZ diagrams are built from: a cell
%  in an array, the one cell under the cursor, a circular graph node,
%  the arrow between two of them. Every book has the whole set:
%
%      utpcell, utphot, utpdone, utpdim, utpwall        array cells
%      utpnode, utpnodehot, utpnodedone, utpnodedim     graph nodes
%      utparrow, utpptr, utpdash, utpedge               connectors
%      utplabel, utptag, utpbrace                       annotation
%      utpstep, utpstage, utppanel, utpbox              boxes
%
%  A cell is 1.0cm by 0.8cm. A book whose diagrams were drawn to a
%  different grid resizes the cells once, after the input:
%
%      \utpdiagramsize{1.05cm}{0.85cm}
%
%  Larger vocabularies for one kind of drawing live in .shared/diagrams/
%  and are loaded by name, after the input:
%
%      \utpusediagrams{network}         devices, links, packet headers,
%                                       sequence diagrams, state machines
%      \utpusediagrams{process}         business-process flows, documents,
%                                       swimlanes, status chips
%      \utpusediagrams{data-objects}    vectors, lists and data frames
%                                       drawn as cells; flowcharts
%      \utpusediagrams{network,process} several at once

\newcommand{\utpdiagramsize}[2]{%
  \tikzset{
    utpcell/.style={draw=UTPBlue,fill=UTPPaleBlue,minimum width=#1,minimum height=#2,
                    font=\small\bfseries\color{UTPNavy},inner sep=1pt},
    utphot/.style={draw=UTPOrange,fill=UTPPaleOrange,minimum width=#1,minimum height=#2,
                   font=\small\bfseries\color{UTPNavy},line width=1pt,inner sep=1pt},
    utpdone/.style={draw=UTPGreen,fill=UTPPaleTeal,minimum width=#1,minimum height=#2,
                    font=\small\bfseries\color{UTPNavy},inner sep=1pt},
    utpdim/.style={draw=UTPMuted!45,fill=UTPPaleGray,minimum width=#1,minimum height=#2,
                   font=\small\color{UTPMuted},inner sep=1pt},
    utpwall/.style={draw=UTPNavy,fill=UTPNavy!82,minimum width=#1,minimum height=#2,
                    inner sep=1pt},
    utpnode/.style={circle,draw=UTPBlue,fill=UTPPaleBlue,minimum size=8mm,
                    font=\small\bfseries\color{UTPNavy},inner sep=0.5pt},
    utpnodehot/.style={circle,draw=UTPOrange,fill=UTPPaleOrange,minimum size=8mm,line width=1pt,
                       font=\small\bfseries\color{UTPNavy},inner sep=0.5pt},
    utpnodedone/.style={circle,draw=UTPGreen,fill=UTPPaleTeal,minimum size=8mm,
                        font=\small\bfseries\color{UTPNavy},inner sep=0.5pt},
    utpnodedim/.style={circle,draw=UTPMuted!45,fill=UTPPaleGray,minimum size=8mm,
                       font=\small\color{UTPMuted},inner sep=0.5pt},
    utparrow/.style={-{Stealth[length=2.2mm]},thick,draw=UTPTeal},
    utpptr/.style={-{Stealth[length=2mm]},semithick,draw=UTPBlue},
    utpdash/.style={-{Stealth[length=2mm]},semithick,draw=UTPRed,dashed},
    utpedge/.style={thick,draw=UTPTeal},
    utplabel/.style={font=\scriptsize\color{UTPMuted},inner sep=1pt},
    utptag/.style={font=\scriptsize\bfseries\color{UTPOrange},inner sep=1pt},
    utpbrace/.style={decorate,decoration={brace,amplitude=4pt,raise=2pt},draw=UTPMuted,semithick},
    utpstep/.style={draw=UTPNavy,fill=white,rounded corners=1.5mm,minimum height=0.75cm,
                    font=\scriptsize\color{UTPNavy},align=center,inner sep=2pt},
    utpstage/.style={draw=UTPNavy,fill=white,rounded corners=1.5mm,minimum height=0.75cm,
                     font=\scriptsize\color{UTPNavy},align=center,inner sep=2pt},
    utppanel/.style={draw=UTPBlue!30,fill=white,rounded corners=1.5mm,inner sep=2.5mm},
    utpbox/.style={draw=UTPBlue,fill=UTPPaleBlue,rounded corners=2mm,align=center,
                   font=\small\color{UTPNavy},inner sep=2.2mm}
  }%
}
\utpdiagramsize{1.0cm}{0.8cm}

% \@for, not \foreach: \foreach runs each pass in a group, which would
% throw away the styles a library declares.
\makeatletter
\newcommand{\utpusediagrams}[1]{%
  \@for\utp@library:=#1\do{\input{../.shared/diagrams/\utp@library}}%
}
\makeatother
